\documentclass[journal]{IEEEtran}

\usepackage{amsmath}
\usepackage{amssymb}
\usepackage{bm}
\usepackage{booktabs}
\usepackage{array}
\usepackage{float}

\begin{document}

\section{Lagrangian Formulation}

The dynamics of the two-wheeled inverted pendulum are derived using the
Euler--Lagrange formulation. The Lagrangian is defined as
\begin{equation}
    \mathcal{L} = T - V,
    \label{eq:lagrangian}
\end{equation}
where $T$ is the total kinetic energy and $V$ is the gravitational
potential energy.

For a system with generalized coordinates $q_i$, the
Euler--Lagrange equations are
\begin{equation}
    \frac{\mathrm{d}}{\mathrm{d}t}
    \left(
        \frac{\partial \mathcal{L}}
        {\partial \dot{q}_i}
    \right)
    -
    \frac{\partial \mathcal{L}}
    {\partial q_i}
    =
    Q_i,
    \label{eq:euler_lagrange}
\end{equation}
where $Q_i$ is the generalized nonconservative force associated with
$q_i$. Because the no-lateral-slip condition is nonholonomic, the
equations for $x$ and $\psi$ require additional terms; these are derived
in the equations-of-motion subsection.

\subsection{Coordinate System}

The system is described using an inertial Cartesian coordinate frame
$\mathcal{F}_I$ with axes $(X,Y,Z)$, where $Z$ is vertical and positive
upward. A heading frame $\mathcal{F}_H$ has its origin at the midpoint of
the wheel axle. Its $x_H$ axis points forward along the robot heading, its
$y_H$ axis points to the left, and its $z_H$ axis is parallel to $Z$.

The position of the axle midpoint is
\begin{equation}
    \bm{r}_{\mathrm{a}}
    =
    \begin{bmatrix}
        X \\
        Y \\
        z
    \end{bmatrix},
    \label{eq:axle_position}
\end{equation}
where $z$ is constant. The heading angle $\psi$ is the rotation of
$\mathcal{F}_H$ about $Z$, positive counterclockwise when viewed from above
(a left turn). The distance traveled along the heading is $x$, defined by
\begin{equation}
    \dot{x} = \dot{X}\cos\psi + \dot{Y}\sin\psi .
    \label{eq:x_definition}
\end{equation}
Note that $x$ is a path length and not an inertial Cartesian coordinate.

The center of mass of the pendulum body is located at
\begin{equation}
    \bm{r}_{\mathrm{b}}
    =
    \bm{r}_{\mathrm{a}}
    +
    \ell
    \begin{bmatrix}
        \sin\theta\cos\psi \\
        \sin\theta\sin\psi \\
        \cos\theta
    \end{bmatrix},
    \label{eq:body_position}
\end{equation}
where $\ell$ is the distance from the wheel axle to the center of mass
of the pendulum body.

The pitch angle $\theta$ is measured clockwise from the positive $Z$ axis
when viewed with $x_H$ pointing to the right, i.e., positive $\theta$ tilts
the body forward. Therefore,
\begin{equation}
    \theta = 0
\end{equation}
corresponds to the upright configuration of the inverted pendulum.

The generalized coordinates are
\begin{equation}
    q_1 = x,
    \quad
    q_2 = \theta,
    \quad
    q_3 = \psi,
    \quad
    q_4 = \phi_{\mathrm l},
    \quad
    q_5 = \phi_{\mathrm r},
    \label{eq:generalized_coordinates}
\end{equation}
where $\phi_{\mathrm l}$ and $\phi_{\mathrm r}$ are the \emph{absolute}
(inertial-frame) angular positions of the left and right wheels, positive
for forward rolling. If wheel angles are measured relative to the body
(e.g., by encoders), then $\phi=\phi^{\mathrm{rel}}+\theta$.

The wheel angular velocities are
\begin{equation}
    \omega_{\mathrm l}
    =
    \dot{\phi}_{\mathrm l},
    \qquad
    \omega_{\mathrm r}
    =
    \dot{\phi}_{\mathrm r}.
    \label{eq:wheel_angular_velocity}
\end{equation}

\subsection{Rolling and Nonholonomic Constraints}

Assuming pure rolling without slipping, the wheel velocities satisfy
\begin{equation}
    \dot{x}
    =
    \frac{r_{\mathrm w}}{2}
    \left(\dot{\phi}_{\mathrm r}+\dot{\phi}_{\mathrm l}\right),
    \qquad
    \dot{\psi}
    =
    \frac{r_{\mathrm w}}{w_{\mathrm{axis}}}
    \left(\dot{\phi}_{\mathrm r}-\dot{\phi}_{\mathrm l}\right),
    \label{eq:rolling_constraint_velocity}
\end{equation}
where $r_{\mathrm w}$ is the wheel radius and $w_{\mathrm{axis}}$ is the
distance between the left and right wheels. The corresponding position
constraints are
\begin{equation}
    x-\frac{r_{\mathrm w}}{2}\left(\phi_{\mathrm r}+\phi_{\mathrm l}\right)=C_x,
    \qquad
    \psi-\frac{r_{\mathrm w}}{w_{\mathrm{axis}}}\left(\phi_{\mathrm r}-\phi_{\mathrm l}\right)=C_\psi,
    \label{eq:rolling_constraint_position}
\end{equation}
where $C_x$ and $C_\psi$ are constants determined by the initial
configuration.

Because the rolling constraints eliminate two degrees of freedom,
$\phi_{\mathrm l}$ and $\phi_{\mathrm r}$ may be eliminated from the final
equations of motion using
\begin{equation}
    \dot{\phi}_{\mathrm{r,l}}
    =
    \frac{1}{r_{\mathrm w}}
    \left(
        \dot{x}
        \pm
        \frac{w_{\mathrm{axis}}}{2}\dot{\psi}
    \right).
    \label{eq:rolling_constraint_acceleration}
\end{equation}
In addition, the wheels cannot slip sideways, so the lateral velocity $u$ of
the axle midpoint in $\mathcal{F}_H$ satisfies the nonholonomic constraint
\begin{equation}
    u = -\dot{X}\sin\psi + \dot{Y}\cos\psi = 0 .
    \label{eq:nonholonomic}
\end{equation}

\subsection{Kinetic Energy}

Let $u$ temporarily denote the lateral velocity of the axle midpoint; it is
set to zero after differentiation. The velocity of the pendulum-body center
of mass, expressed in $\mathcal{F}_H$, has components
\begin{equation}
    \bm{v}_{\mathrm b}
    =
    \begin{bmatrix}
        \dot{x} + \ell\dot{\theta}\cos\theta \\
        u + \ell\dot{\psi}\sin\theta \\
        -\ell\dot{\theta}\sin\theta
    \end{bmatrix}.
    \label{eq:body_velocity}
\end{equation}
Setting $u=0$ and using $\dot z=0$, the translational kinetic energy of the
wheel assembly and pendulum body is
\begin{equation}
    T_{\mathrm{trans}}
    =
    \frac{1}{2}
    (m_{\mathrm b}+m_{\mathrm w})
    \dot{x}^{\,2},
    \label{eq:translational_kinetic_energy}
\end{equation}
where $m_{\mathrm b}$ is the mass of the pendulum body and $m_{\mathrm w}$ is
the combined mass of both wheels. The additional translational energy of the
pendulum-body center of mass due to its angular displacement is
\begin{equation}
    T_{\mathrm{corr}}
    =
    m_{\mathrm b}\ell
    \dot{x}\dot{\theta}\cos\theta
    +
    \frac{1}{2}
    m_{\mathrm b}\ell^2
    \dot{\theta}^{\,2}.
    \label{eq:correction_kinetic_energy}
\end{equation}

The pitch rotational kinetic energy of the pendulum body is
\begin{equation}
    T_{\mathrm b,\theta}
    =
    \frac{1}{2}
    I_{\mathrm b:2}
    \dot{\theta}^{\,2}.
    \label{eq:body_rotational_energy}
\end{equation}
The yaw kinetic energy of the pendulum body is
\begin{equation}
    T_{\mathrm b,\psi}
    =
    \frac{1}{2}
    \dot{\psi}^{\,2}
    \left(
        \sin^2\theta
        \left(I_{\mathrm{b}:1}+m_{\mathrm b}\ell^2\right)
        +\cos^2\theta\, I_{\mathrm{b}:3}
    \right),
    \label{eq:body_yaw_energy}
\end{equation}
where the $m_{\mathrm b}\ell^2\sin^2\theta$ term arises from the lateral
motion of the center of mass during yaw. The wheels do not pitch with
$\theta$, so their yaw inertia is constant:
\begin{equation}
    T_{\mathrm w,\psi}
    =
    \frac{1}{2}
    \dot{\psi}^{\,2}
    \left(
        2I_{\mathrm{w,d}}
        +
        \frac{m_{\mathrm w}w_{\mathrm{axis}}^{\,2}}{4}
    \right),
    \label{eq:wheel_yaw_energy}
\end{equation}
where $I_{\mathrm{w,d}}$ is the inertia of one wheel about a diameter, and
the second term is the parallel-axis contribution of the two wheels. The
wheel spin energy is
\begin{equation}
    T_{\mathrm w,\mathrm{spin}}
    =
    \frac{1}{2}
    I_{\mathrm w}
    \left(
        \dot{\phi}_{\mathrm l}^{\,2}
        +
        \dot{\phi}_{\mathrm r}^{\,2}
    \right)
    =
    \frac{I_{\mathrm w}}{r_{\mathrm w}^{\,2}}\dot{x}^{\,2}
    +
    \frac{1}{2}
    \frac{I_{\mathrm w}w_{\mathrm{axis}}^{\,2}}{2r_{\mathrm w}^{\,2}}
    \dot{\psi}^{\,2},
    \label{eq:wheel_rotational_energy}
\end{equation}
where $I_{\mathrm w}$ is the spin inertia of a single wheel about its axle
and \eqref{eq:rolling_constraint_acceleration} has been applied.

Define the effective translational mass, the effective pitch inertia, and
the yaw inertia
\begin{align}
    m_{\mathrm{eff}}
    &=
    m_{\mathrm b}+m_{\mathrm w}+\frac{2I_{\mathrm w}}{r_{\mathrm w}^{\,2}},
    \label{eq:m_eff}
    \\
    I_{\theta}
    &=
    I_{\mathrm b:2}+m_{\mathrm b}\ell^2,
    \label{eq:I_theta}
    \\
    \mathcal{J}(\theta)
    &=
    J_0 + \Delta I\,\sin^2\theta,
    \label{eq:yaw_inertia}
\end{align}
with
\begin{align}
    J_{\mathrm w}
    &=
    2I_{\mathrm{w,d}}
    +
    \frac{m_{\mathrm w}w_{\mathrm{axis}}^{\,2}}{4}
    +
    \frac{I_{\mathrm w}w_{\mathrm{axis}}^{\,2}}{2r_{\mathrm w}^{\,2}},
    \nonumber\\
    J_0
    &=
    J_{\mathrm w} + I_{\mathrm b:3},
    \nonumber\\
    \Delta I
    &=
    I_{\mathrm b:1} + m_{\mathrm b}\ell^2 - I_{\mathrm b:3}.
    \label{eq:yaw_inertia_constants}
\end{align}
Thus $\mathcal{J}(\theta)$ equals the sum of the wheel yaw inertia and
\eqref{eq:body_yaw_energy}.

The total kinetic energy, after applying the rolling constraints, is
\begin{align}
    T
    &=
    T_{\mathrm{trans}}
    +
    T_{\mathrm{corr}}
    +
    T_{\mathrm b,\theta}
    +
    T_{\mathrm b,\psi}
    +
    T_{\mathrm w,\psi}
    +
    T_{\mathrm w,\mathrm{spin}}
    \nonumber\\
    &=
    \frac{1}{2}
    m_{\mathrm{eff}}
    \dot{x}^{\,2}
    +
    m_{\mathrm b}\ell
    \dot{x}\dot{\theta}\cos\theta
    \nonumber\\
    &\quad
    +
    \frac{1}{2}
    I_{\theta}
    \dot{\theta}^{\,2}
    +
    \frac{1}{2}
    \mathcal{J}(\theta)
    \dot{\psi}^{\,2}.
    \label{eq:constrained_kinetic_energy}
\end{align}

When $u\neq 0$ is retained, the only additional term that survives in the
derivatives evaluated at $u=0$ comes from the lateral kinetic energy
$\tfrac12 m_{\mathrm b}\left(u+\ell\dot\psi\sin\theta\right)^2$ of the
pendulum body, giving
\begin{equation}
    \left.\frac{\partial T}{\partial u}\right|_{u=0}
    =
    m_{\mathrm b}\ell\,\dot{\psi}\sin\theta .
    \label{eq:dT_du}
\end{equation}

\subsection{Potential Energy}

The gravitational potential energy of the pendulum body is
\begin{equation}
    V
    =
    m_{\mathrm b}g Z_{\mathrm b}
    =
    m_{\mathrm b}g
    \left(
        z+\ell\cos\theta
    \right).
    \label{eq:potential_full}
\end{equation}

Since $z$ is constant, the term $m_{\mathrm b}gz$ is constant and does
not affect the equations of motion. It can therefore be omitted,
yielding
\begin{equation}
    V
    =
    m_{\mathrm b}g\ell\cos\theta.
    \label{eq:potential_energy}
\end{equation}

\subsection{Constrained Lagrangian}

Using \eqref{eq:constrained_kinetic_energy} and
\eqref{eq:potential_energy}, the constrained Lagrangian is
\begin{align}
    \mathcal{L}
    &=
    \frac{1}{2}
    m_{\mathrm{eff}}
    \dot{x}^{\,2}
    +
    m_{\mathrm b}\ell
    \dot{x}\dot{\theta}\cos\theta
    \nonumber\\
    &\quad
    +
    \frac{1}{2}
    I_{\theta}
    \dot{\theta}^{\,2}
    +
    \frac{1}{2}
    \mathcal{J}(\theta)
    \dot{\psi}^{\,2}
    -
    m_{\mathrm b}g\ell\cos\theta.
    \label{eq:constrained_lagrangian}
\end{align}

\subsection{Generalized Forces}

Let $M_{\mathrm{lw}}$ and $M_{\mathrm{rw}}$ denote the torques produced by
the left and right motors. Each motor applies a torque $M$ to the pendulum
body in the positive $\theta$ direction and an equal and opposite torque
$-M$ to its wheel. The virtual work is
\begin{equation}
    \delta W
    =
    \left(M_{\mathrm{lw}}+M_{\mathrm{rw}}\right)\delta\theta
    -
    M_{\mathrm{lw}}\,\delta\phi_{\mathrm l}
    -
    M_{\mathrm{rw}}\,\delta\phi_{\mathrm r}.
    \label{eq:virtual_work}
\end{equation}
Substituting the constraints
$\delta\phi_{\mathrm{r,l}}=\left(\delta x\pm\tfrac{w_{\mathrm{axis}}}{2}\delta\psi\right)/r_{\mathrm w}$
gives the generalized forces
\begin{align}
    Q_{\theta}
    &=
    M_{\mathrm{lw}}
    + M_{\mathrm{rw}},
    \label{eq:pendulum_motor_torque}
    \\
    Q_x
    &=
    -\frac{M_{\mathrm{lw}} + M_{\mathrm{rw}}}{r_{\mathrm w}},
    \label{eq:effective_motor_force}
    \\
    Q_{\psi}
    &=
    \frac{w_{\mathrm{axis}}}{2r_{\mathrm w}}
    \left(
        M_{\mathrm{lw}}
        -
        M_{\mathrm{rw}}
    \right).
    \label{eq:yaw_torque}
\end{align}
Positive $M$ tilts the body forward ($+\theta$) and drives the vehicle
backward. A larger right-wheel speed turns the robot to the left, which is
the positive $\psi$ direction. The yaw torque sign follows from
$\dot\psi = r_{\mathrm w}(\dot\phi_{\mathrm r}-\dot\phi_{\mathrm l})/w_{\mathrm{axis}}$.

\subsection{Equations of Motion}

Because $x$ and $\psi$ are subject to the nonholonomic constraint
\eqref{eq:nonholonomic}, the equations for these two coordinates are
obtained from the Hamel (quasi-velocity) form of Lagrange's equations,
\begin{align}
    \frac{\mathrm{d}}{\mathrm{d}t}
    \frac{\partial \mathcal{L}}{\partial \dot{x}}
    -
    \dot{\psi}
    \left.\frac{\partial \mathcal{L}}{\partial u}\right|_{u=0}
    &=
    Q_x,
    \label{eq:hamel_x}
    \\
    \frac{\mathrm{d}}{\mathrm{d}t}
    \frac{\partial \mathcal{L}}{\partial \dot{\psi}}
    +
    \dot{x}
    \left.\frac{\partial \mathcal{L}}{\partial u}\right|_{u=0}
    &=
    Q_\psi,
    \label{eq:hamel_psi}
\end{align}
while $\theta$ is an ordinary holonomic coordinate and uses
\eqref{eq:euler_lagrange} directly.

\subsubsection{Translational Equation}

The momentum conjugate to $\dot x$ is
\begin{equation}
    \frac{\partial \mathcal{L}}
    {\partial \dot{x}}
    =
    m_{\mathrm{eff}}
    \dot{x}
    +
    m_{\mathrm b}\ell
    \dot{\theta}\cos\theta.
    \label{eq:partial_L_partial_xdot}
\end{equation}

Taking the time derivative gives
\begin{equation}
    \frac{\mathrm{d}}{\mathrm{d}t}
    \left(
        \frac{\partial \mathcal{L}}
        {\partial \dot{x}}
    \right)
    =
    m_{\mathrm{eff}}
    \ddot{x}
    +
    m_{\mathrm b}\ell
    \left(
        \ddot{\theta}\cos\theta
        -
        \dot{\theta}^{\,2}\sin\theta
    \right).
    \label{eq:d_dt_partial_L_partial_xdot}
\end{equation}

Using \eqref{eq:dT_du} in \eqref{eq:hamel_x}, the translational equation of
motion is
\begin{align}
    &m_{\mathrm{eff}}
    \ddot{x}
    +
    m_{\mathrm b}\ell
    \cos\theta\,\ddot{\theta}
    -
    m_{\mathrm b}\ell
    \sin\theta
    \left(
        \dot{\theta}^{\,2}
        +
        \dot{\psi}^{\,2}
    \right)
    \nonumber\\
    &\qquad
    =
    -\frac{M_{\mathrm{lw}}+M_{\mathrm{rw}}}{r_{\mathrm w}}.
    \label{eq:translation_eom}
\end{align}
The $\dot\psi^2$ term is the centripetal term arising from the lateral offset
of the center of mass while turning.

\subsubsection{Pendulum Angular Equation}

Applying the Euler--Lagrange equation to $\theta$ gives
\begin{equation}
    \frac{\partial \mathcal{L}}
    {\partial \dot{\theta}}
    =
    m_{\mathrm b}\ell
    \dot{x}\cos\theta
    +
    I_{\theta}
    \dot{\theta}.
    \label{eq:partial_L_partial_thetadot}
\end{equation}

Taking the time derivative gives
\begin{equation}
    \frac{\mathrm{d}}{\mathrm{d}t}
    \left(
        \frac{\partial \mathcal{L}}
        {\partial \dot{\theta}}
    \right)
    =
    m_{\mathrm b}\ell
    \left(
        \ddot{x}\cos\theta
        -
        \dot{x}\dot{\theta}\sin\theta
    \right)
    +
    I_{\theta}
    \ddot{\theta}.
    \label{eq:d_dt_partial_L_partial_thetadot}
\end{equation}

The partial derivative with respect to $\theta$ is
\begin{align}
    \frac{\partial \mathcal{L}}
    {\partial \theta}
    &=
    -m_{\mathrm b}\ell
    \dot{x}\dot{\theta}\sin\theta
    +
    \frac{1}{2}\dot{\psi}^{\,2}\,
    \Delta I\,\sin2\theta
    \nonumber\\
    &\quad
    +
    m_{\mathrm b}g\ell\sin\theta.
    \label{eq:partial_L_partial_theta}
\end{align}

Therefore, the pitch ($\theta$) equation of motion is
\begin{align}
    &I_{\theta}
    \ddot{\theta}
    +
    m_{\mathrm b}\ell
    \cos\theta\,\ddot{x}
    -
    \frac{1}{2}\dot{\psi}^{\,2}\,
    \Delta I\,\sin2\theta
    -
    m_{\mathrm b}g\ell\sin\theta
    \nonumber\\
    &\qquad
    =
    M_{\mathrm{lw}} + M_{\mathrm{rw}}.
    \label{eq:angular_eom}
\end{align}
The first term represents the rotational inertia of the pendulum body, the
second represents translational--rotational coupling, the third represents
yaw-induced torque, and the fourth represents the gravitational torque.

\subsubsection{Yaw Equation of Motion}

The momentum conjugate to $\dot\psi$ is
\begin{equation}
    \frac{\partial \mathcal{L}}{\partial \dot{\psi}}
    =
    \mathcal{J}(\theta)\,\dot{\psi}.
\end{equation}
Taking the time derivative and using $\mathrm{d}\mathcal{J}/\mathrm{d}t=\Delta I\sin2\theta\,\dot\theta$ gives
\begin{equation}
    \frac{\mathrm{d}}{\mathrm{d}t}\frac{\partial \mathcal{L}}{\partial \dot{\psi}}
    =
    \mathcal{J}(\theta)\,\ddot{\psi}
    +
    \Delta I\,\sin2\theta\,\dot{\theta}\dot{\psi}.
\end{equation}
Using \eqref{eq:dT_du} in \eqref{eq:hamel_psi}, the yaw equation of motion is
\begin{equation}
    \mathcal{J}(\theta)\,\ddot{\psi}
    +
    \Delta I\,\sin2\theta\,\dot{\theta}\dot{\psi}
    +
    m_{\mathrm b}\ell\sin\theta\,\dot{x}\dot{\psi}
    =
    \frac{w_{\mathrm{axis}}}{2r_{\mathrm w}}
    \left(
        M_{\mathrm{lw}}
        -
        M_{\mathrm{rw}}
    \right).
    \label{eq:yaw_eom}
\end{equation}
The second term is the gyroscopic coupling between pitch rate and yaw rate,
and the third term comes from the nonholonomic constraint.

\subsection{Symbol and Coordinate Definitions}

\begin{table}[H]
\centering
\caption{Symbols and coordinate definitions.}
\label{tab:symbol_definitions}
\renewcommand{\arraystretch}{1.1}
\begin{tabular}{
    >{\raggedright\arraybackslash}p{0.22\linewidth}
    >{\raggedright\arraybackslash}p{0.68\linewidth}
}
\toprule
Symbol & Definition \\
\midrule

$\mathcal{F}_I$
& Inertial Cartesian frame with axes $(X,Y,Z)$; $Z$ positive upward. \\

$\mathcal{F}_H$
& Heading frame at the axle midpoint: $x_H$ forward, $y_H$ left,
$z_H$ vertical. \\

$X,Y$
& Inertial horizontal position of the axle midpoint. \\

$z$
& Height of the axle above the ground; constant. \\

$x$
& Distance traveled along the heading, $\dot{x}=\dot{X}\cos\psi+\dot{Y}\sin\psi$
(path length, not an inertial coordinate). \\

$u$
& Lateral velocity of the axle midpoint in $\mathcal{F}_H$; constrained to
zero (no lateral slip), used only to form the nonholonomic terms. \\

$\psi$
& Heading (yaw) angle about $Z$, positive counterclockwise viewed from
above (left turn). \\

$\theta$
& Pitch angle of the pendulum body from the positive vertical axis, positive
forward; $\theta=0$ is upright. \\

$\phi_{\mathrm l},\phi_{\mathrm r}$
& Absolute (inertial-frame) angular positions of the left and right wheels,
positive for forward rolling. \\

$\omega_{\mathrm l},\omega_{\mathrm r}$
& Wheel angular velocities, $\omega_{\mathrm{l,r}}=\dot{\phi}_{\mathrm{l,r}}$. \\

$\bm r_{\mathrm a}$
& Position vector of the axle midpoint. \\

$\bm r_{\mathrm b}$
& Position vector of the pendulum-body center of mass. \\

$Z_{\mathrm b}$
& Vertical coordinate of the pendulum-body center of mass. \\

$\bm v_{\mathrm b}$
& Velocity of the pendulum-body center of mass in $\mathcal{F}_H$. \\

$\ell$
& Distance from the wheel axle to the pendulum-body center of mass. \\

$r_{\mathrm w}$
& Wheel radius. \\

$w_{\mathrm{axis}}$
& Distance between the left and right wheels (track width). \\

$m_{\mathrm b}$
& Mass of the pendulum body. \\

$m_{\mathrm w}$
& Combined mass of both wheels (each wheel center of mass on the axle). \\

$I_{\mathrm b:1}, I_{\mathrm b:2}, I_{\mathrm b:3}$
& Principal moments of inertia of the pendulum body about its center of
mass: axis 1 forward, axis 2 the pitch axis (perpendicular to the plane of
motion), axis 3 along the body's upright axis. \\

$I_{\mathrm w}$
& Spin moment of inertia of a single wheel about its axle. \\

$I_{\mathrm{w,d}}$
& Moment of inertia of a single wheel about a diameter (used for yaw). \\

$m_{\mathrm{eff}}$
& Effective translational mass,
$m_{\mathrm b}+m_{\mathrm w}+2I_{\mathrm w}/r_{\mathrm w}^2$. \\

$I_\theta$
& Effective pitch inertia, $I_{\mathrm b:2}+m_{\mathrm b}\ell^2$. \\

$J_{\mathrm w}$
& Constant wheel yaw inertia,
$2I_{\mathrm{w,d}}+m_{\mathrm w}w_{\mathrm{axis}}^2/4+I_{\mathrm w}w_{\mathrm{axis}}^2/(2r_{\mathrm w}^2)$. \\

$J_0$
& Yaw inertia at $\theta=0$, $J_{\mathrm w}+I_{\mathrm b:3}$. \\

$\Delta I$
& $I_{\mathrm b:1}+m_{\mathrm b}\ell^2-I_{\mathrm b:3}$. \\

$\mathcal{J}(\theta)$
& Total yaw inertia, $J_0+\Delta I\sin^2\theta$. \\

$g$
& Gravitational acceleration. \\

$T$
& Total kinetic energy of the system. \\

$V$
& Gravitational potential energy of the system. \\

$\mathcal{L}$
& Lagrangian, $\mathcal{L}=T-V$. \\

$q_i$
& $i$th generalized coordinate. \\

$Q_i$
& Generalized nonconservative force associated with $q_i$. \\

$M_{\mathrm{lw}},M_{\mathrm{rw}}$
& Left and right motor torques; each acts on the body in the positive
$\theta$ direction and on its wheel in the opposite direction. \\

$Q_x,Q_\theta,Q_\psi$
& Generalized forces for $x$, $\theta$, and $\psi$ from the motor torques. \\

$C_x,C_\psi$
& Constants determined by the initial rolling configuration. \\

$t$
& Time. \\

\bottomrule
\end{tabular}
\end{table}

\end{document}